\documentclass[11pt,a4paper]{article}
\usepackage[margin=1in]{geometry}
\usepackage{booktabs}
\usepackage{longtable}
\usepackage{graphicx}
\usepackage{hyperref}
\usepackage{listings}
\usepackage{xcolor}
\usepackage{enumitem}
\usepackage{caption}
\usepackage{amsmath}
\usepackage{amssymb}
\usepackage{tcolorbox}
\tcbuselibrary{breakable}

\hypersetup{
    colorlinks=true,
    linkcolor=blue,
    filecolor=magenta,
    urlcolor=cyan,
}

\tcbset{
    promptbox/.style={
        colback=blue!5!white,
        colframe=blue!75!black,
        title={#1},
        fonttitle=\bfseries,
        before skip=6pt,
        after skip=6pt,
        breakable,
    }
}

\title{Performance Report 21: Vulkan Engine \\
\large Solid and Water Rasterization --- Double Geometry, Always-On Tessellation, and the Fragment Fetch Multiplier}
\author{}
\date{\today}

\begin{document}

\maketitle

\begin{abstract}
Reports~17--20 audited the hybrid ray-tracing paths (17, 18), the frame-level
water machinery (19), and the water noise constant factor (20). The
\emph{raster} core --- the solid terrain forward/deferred passes and the water
geometry pass that every preset executes --- has never had a dedicated report.
This report closes that gap with a static audit of the shipped solid and water
rasterization code plus SPIR-V measurements of the built modules.

It finds \textbf{15 issues}: the solid pass has no non-tessellation pipeline
family, so the default configuration (tessellation off) still runs every solid
triangle through the TCS, the fixed-function tessellator and the TES (C1); the
deferred depth$+$color scheme draws the full tessellated mesh \emph{twice},
and the depth pass pays the full height-map displacement sampling even though
its fragment shader is 240 words long (C2); triplanar shading multiplies every
material map fetch by up to $3\times$ projections $\times$ $3\times$ materials
(36 fetches per pixel worst case), confirmed by 34 static sample sites in
\texttt{main.frag.spv} (C3). Behind them: a 3-cascade shadow pass that also
runs the full tessellation$+$displacement pipeline plus six fullscreen EVSM
blurs (H4), an unconditional per-pixel RGB$\rightarrow$HSV$\rightarrow$RGB
round-trip with divergent branches (H5), secondary CSM lookups from inside RT
hit shading (H6), a dead 12-byte-per-vertex color attribute carried through
four stages (H7), and a TCS that computes every edge level three times (H8).
The water findings are residual by design (reports~19--20 already landed):
implicit-LOD taps in divergent composite flow (M9), a \texttt{waterRenderScale}
that defaults to 1.0 and that neither preset touches (M10), a second full
water cull dispatch whenever shadows are on (M11), and vertex-texture-fetch
shore measurement in the no-tess water vertex shader (M12).

Each issue is paired with the standard technique that fixes it and a
copy-pasteable prompt (goal, steps, files, validation) preserving
validation-clean execution, Synchronization2, runtime feature detection, and
the raster$+$CSM fallback required by \texttt{AGENTS.md}. The application
scores \textbf{6.9/10} overall (\S\ref{sec:score}): a strong submission and
culling architecture held back by geometry that is processed two to three
times per frame and a fragment stage that fetches the same texel families
repeatedly instead of once.
\end{abstract}

\tableofcontents
\newpage

% =====================================================================
\section{Executive Summary}
% =====================================================================

\subsection{Scope and Method}

This report covers \emph{only} the rasterization-time cost of the two surfaces
that exist in every frame: solid terrain
(\texttt{SolidRenderer}, \texttt{main.vert}/\texttt{main.tesc}/\texttt{main.tese}/\texttt{main.frag},
\texttt{includes/solid\_surface.glsl}, the CSM shadow path) and water
(\texttt{WaterRenderer}, the water variants of the same stages, the water
geometry pass and the \texttt{postprocess.frag} composite). Ray-tracing paths
(reports~17--18), water frame machinery (report~19) and water noise ALU
(report~20) are referenced but not re-audited; where this report touches them
it is only at the interface (e.g.\ secondary CSM lookups from hit shading).

Method: full reads of the renderer sources and shader includes cited below,
plus measurements of the shipped \texttt{bin/shaders/*.spv} modules with
\texttt{spirv-dis} (disassembled-operation counts and static
\texttt{OpImageSample} sites). The application was \emph{not} executed: no GPU
timestamps are quoted, and every dynamic magnitude is an analytic estimate
from verified code structure, clearly labelled as such. No finding is reported
without a file, a line number and a code quote.

\subsection{What Previous Reports Left on the Table}

\begin{itemize}
    \item Report~19's C1 built a \textbf{non-tessellation pipeline family for
    water} and report~20 confirms all twelve report~19 items landed. The
    \textbf{solid} pass never received the same treatment: every solid
    pipeline unconditionally includes the TCS/TCS/TES stages (C1).
    \item The deferred depth$+$color solid scheme (depth prepass $\rightarrow$
    sky $\rightarrow$ color with loaded depth) is \emph{correct} --- the color
    pass enjoys near-perfect early-$z$ rejection --- but its price is paid in
    the vertex/tessellation front-end, which runs twice, displacement
    included (C2).
    \item Report~20 measured fragment \emph{ALU} (46k-instruction water
    module). This report measures fragment \emph{fetch multiplicity} on the
    solid side: the ALU was gated, the fetches were not (C3).
\end{itemize}

\subsection{Measured Module Budget}

\texttt{spirv-dis} on the shipped modules (exact flags as the Makefile,
\texttt{glslc --target-env=vulkan1.3 -O}):

\begin{table}[h]
\centering
\small
\begin{tabular}{@{}lrrrl@{}}
\toprule
\textbf{Module} & \textbf{Words} & \textbf{Disasm.\ ops} & \textbf{Static image samples} & \textbf{Role} \\
\midrule
\texttt{main.frag.spv} & 22\,726 & 5\,394 & \textbf{34} & solid surface (non-RT) \\
\texttt{main.vert.spv} & 3\,232 & --- & --- & solid VS (trivial) \\
\texttt{main.tesc.spv} & 5\,233 & --- & --- & solid TCS \\
\texttt{main.tese.spv} & 7\,881 & --- & --- & solid TES $+$ displacement \\
\texttt{depth\_only.frag.spv} & 240 & --- & 0 & solid depth prepass FS \\
\texttt{main\_water.frag.spv} & 27\,876 & 5\,750 & 12 & water surface (non-RT) \\
\texttt{main\_water\_no\_tess.vert.spv} & 4\,038 & 858 & --- & water no-tess VS \\
\bottomrule
\end{tabular}
\caption{Shipped module sizes. The solid depth-prepass FS is 240 words
($\sim$1\% of the solid FS) but its pipeline runs the same 13k-word
VS$+$TCS$+$TES front-end as the color pass (C2).}
\end{table}

\subsection{Recommended Order}

Port the water C1 treatment to solid (C1) $\rightarrow$ stop displacing the
depth prepass (C2) $\rightarrow$ collapse the triplanar fetch fan-out (C3)
$\rightarrow$ cheapen the shadow path (H4) $\rightarrow$ gate the HSV
round-trip (H5). The water items (M9--M12) and hygiene items (L13--L15) are
independent and cheap.

% =====================================================================
\section{Cost Model}
% =====================================================================

Per-frame raster cost of the two surfaces:

\begin{equation}
C = \underbrace{G_{\mathrm{passes}} \times V \times C_{\mathrm{vert}}}_{\text{geometry, \S3}}
  + \underbrace{\sum_{\mathrm{px}} F_{\mathrm{fetch}} \times C_{\mathrm{tex}}}_{\text{fragment fetches, \S3}}
  + \underbrace{\sum_{\mathrm{px}} A_{\mathrm{alu}}}_{\text{fragment ALU (report~20)}}
  + \underbrace{C_{\mathrm{shadow}} + C_{\mathrm{composite}} + C_{\mathrm{cpu}}}_{\text{\S4--\S5}}
\end{equation}

with $G_{\mathrm{passes}} \approx 2$ (solid, deferred) $+ 3$ (shadow cascades,
tessellated) $+ 1$--$2$ (water geometry, plus back-face when volume is
needed). The dominant, least-gated term is $G_{\mathrm{passes}}$: each pass
re-runs vertex fetch, tessellation and (for solid) height-map displacement.
The second is $F_{\mathrm{fetch}}$ on solid pixels, where triplanar projection
turns 4 map reads into up to 36.

% =====================================================================
\section{Critical Issues}
% =====================================================================

\subsection{C1: Solid Has No Non-Tessellation Pipeline Family --- Tessellation Off Still Pays TCS $+$ Tessellator $+$ TES}
\label{sec:c1}

\textbf{Severity: Critical.}

Report~19 C1 built a \texttt{TRIANGLE\_LIST}, TCS/TES-free pipeline family for
water and measured the win. The solid pass never got one. Every solid pipeline
in \texttt{SolidRenderer::createPipelines} unconditionally chains the
tessellation stages:

\begin{lstlisting}[basicstyle=\footnotesize\ttfamily]
auto [pipeline, layout] = app->createGraphicsPipeline(
    { vertexShader.info, tescShader.info, teseShader.info,
      fragmentShader.info }, ... );           // SolidRenderer.cpp:268-280
\end{lstlisting}

and the same four-stage chain is repeated for the RT variant (:287--302), the
profiling variant (:305--321), the depth prepass (:326--340), the deferred
depth (:353--363) and deferred color (:370--378) pipelines, the brush variants
(:422--456) and the wireframe pipeline (\texttt{SolidRenderer.cpp:618--625}).
A \texttt{grep} for \texttt{NoTess}/\texttt{NO\_TESS} in
\texttt{SolidRenderer.*} returns nothing: no non-tessellation solid pipeline
exists.

Meanwhile tessellation defaults to \textbf{off}:
\texttt{utils/Settings.hpp:62} (\texttt{tessellationEnabled = false}) and the
Minimal preset explicitly disables it
(\texttt{utils/GraphicsSettingsCommand.hpp}: \texttt{applyMinimal} sets
\texttt{tessellationEnabled = false}). So the default configuration pushes
every solid triangle through three TCS invocations per patch, the
fixed-function tessellator, and the full TES (7\,881 words, including the
barycentric interpolation of nine varyings in
\texttt{solid\_tese.glsl:46--52}) just to emit tessellation level 1.0. The
water path proves the fix works and the plumbing pattern (variant selection at
bind time) already exists in both renderers.

\textit{Technique: pipeline-family specialization on a static toggle.}
Build \texttt{TRIANGLE\_LIST} solid pipelines without TCS/TES (reusing the
water C1 pattern), select them when \texttt{ubo.tessellationEnabled} is false,
and move the TES-only computations (light-space position, sharp normal) into
the VS for that family. Keep the tessellated family for Maximum.

\begin{tcolorbox}[promptbox={Prompt --- C1: non-tessellation solid pipeline family}]
Goal: bind a TCS/TES-free solid pipeline family whenever tessellation is off,
mirroring report~19 C1 for water. Steps: (1) add a \texttt{WATER\_NO\_TESS}-style
compile path (e.g.\ \texttt{SOLID\_NO\_TESS}) to \texttt{main.vert} that writes
the fragment interface directly (\texttt{fragPosLightSpace},
\texttt{fragSharpNormal} from derivatives or the vertex normal,
\texttt{fragTexWeights} from the vertex brush index); (2) create the matching
pipelines in \texttt{SolidRenderer::createPipelines} for the forward, deferred
depth/color and shadow-entry configurations; (3) select the family at bind time
from the same toggle that already collapses tess levels to 1.0. Files:
\texttt{vulkan/renderer/SolidRenderer.cpp}, \texttt{shaders/main.vert},
\texttt{shaders/includes/solid\_tese.glsl}. Validation: validation-clean run,
screenshots identical with tessellation off except multi-material boundary
triangles (approach A, implemented: flat provoking-vertex material in
\texttt{main\_solid\_no\_tess.vert.spv} --- exact for single-material
triangles, which a VS cannot distinguish on indexed draws; RT/profiling/brush
paths keep the tessellated pipelines as a correct fallback), timestamp queries
showing the solid passes shrinking; Maximum (tessellation on) must be
pixel-identical.
\\ \textbf{Correction after implementation:} approach A was reverted for the
color passes. \texttt{LandBrush::paint} assigns materials per vertex by slope
and height bands (\texttt{utils/LandBrush.cpp}), so multi-material triangles
cover entire mountainsides --- flat-shading them visibly breaks blending far
beyond boundary pixels. The shipped fix keeps the tessellated family for all
color passes (blending intact) and binds the no-tess twin only where the
fragment shader samples no materials and the twin is exactly equivalent: the
deferred depth prepass and the EVSM shadow pass. The exact path for the color
pass remains a geometry-shader compression stage (all three corners visible)
or a mesh split at material seams.
\end{tcolorbox}

\subsection{C2: Deferred Depth $+$ Color Draws the Full Tessellated Mesh Twice --- the Depth Pass Displaces Too}
\label{sec:c2}

\textbf{Severity: Critical.}

\texttt{MyApp.cpp:1790} records \texttt{drawDepth} (depth-only,
\texttt{LESS}, no color attachment) and \texttt{MyApp.cpp:1869} records
\texttt{drawColor} (full \texttt{main.frag}, \texttt{LESS\_OR\_EQUAL}, no depth
write) over the same culled command list. The design intent is sound: the
color pass loads prepass depth and gets near-zero overdraw of the expensive
fragment shader. But both passes run the \emph{entire} VS$+$TCS$+$TES
front-end, and the TES displaces unconditionally:

\begin{lstlisting}[basicstyle=\footnotesize\ttfamily]
vec3 displacedLocalPos = mappingFlag > 0.5
    ? applyDisplacement(localPos, localNormal, ...) : localPos;  // solid_tese.glsl:61
\end{lstlisting}

\texttt{applyDisplacement} samples up to three height maps
(\texttt{displacement.glsl:22--36}, triplanar height sampling per material)
\emph{per tessellated vertex} --- in the depth pass, whose fragment shader is
240 words and writes nothing. Concretely, with tessellation on, the most
expensive vertex work in the engine (tessellation $+$ 3 height fetches per
generated vertex) executes twice per frame; with tessellation off, the
redundant TCS/TES stages of C1 execute twice per frame. The depth pass exists
to save fragment work, but it doubles vertex work that, at high tess levels,
dominates.

\textit{Technique: cheap-depth prepass.} Two composable options: (a) a
depth-only TES variant that skips \texttt{applyDisplacement} (conservative
depth: displace by the per-patch \emph{minimum} height instead, or accept the
undisplaced mesh when displacement amplitude is small --- depth only needs to
be conservative, not exact); (b) skip the prepass entirely when the solid
overdraw estimate is low (tessellation off $+$ shadows off already implies a
cheap fragment stage: no CSM fetches, \S\ref{sec:h6}). Option (a) preserves the
early-$z$ benefit for the color pass at a fraction of the vertex cost.

\begin{tcolorbox}[promptbox={Prompt --- C2: displacement-free (conservative) depth prepass}]
Goal: the depth prepass must never pay full displacement sampling. Steps:
(1) add a \texttt{DEPTH\_PASS} compile path to the solid TES that replaces
\texttt{applyDisplacement} with the undisplaced position (or a per-patch
minimum-height displacement for conservatism); (2) build the deferred depth
pipeline (\texttt{SolidRenderer.cpp:342--363}) with it; (3) optionally gate the
whole prepass on (tessellation on \emph{or} shadows on \emph{or} RT on) so
Minimal runs a single forward pass. Files:
\texttt{shaders/includes/solid\_tese.glsl},
\texttt{vulkan/renderer/SolidRenderer.cpp}, \texttt{MyApp.cpp} ($\sim$1790,
$\sim$1869). Validation: depth-buffer screenshots identical within one depth
quantum, color pass pixel-identical, timestamp queries on query slots 6/7 vs
10/11 showing the depth pass shrinking; validation-clean.
\\ \textbf{Correction after implementation:} sub-item (a) --- an undisplaced
(depth-only) prepass --- was declined as unsound and not implemented. The
color pass tests \texttt{LESS\_OR\_EQUAL} against prepass depth, so prepass
depth must equal color depth per pixel: wherever displacement pushes the
surface \emph{away} from the camera, an undisplaced prepass records a nearer
depth, the true surface fails the test, and sky shows through the terrain. A
prepass is either exact or harmful; no cheap approximation preserves both
correctness and the early-$z$ benefit. What shipped instead is (b) as an
opt-in gate: \texttt{Settings::solidDepthPrepass} (default \texttt{true} =
current behavior) skips Instance~1 and binds a depth-write color twin
(\texttt{LESS} over cleared depth), forced back on while solid RT paths run
(the RT variants have no depth-write twin). Tessellation-off frames already
pay almost nothing for the prepass via the C1 no-tess twin, so the gate's
expected win is the tessellation-on, vertex-bound regime --- confirm via the
query slots before flipping the default.
\end{tcolorbox}

\subsection{C3: Triplanar Shading Multiplies Every Map Fetch by up to $9\times$ --- 36 Fetches per Pixel Worst Case}
\label{sec:c3}

\textbf{Severity: Critical.}

The solid fragment shader samples four map families (albedo, normal,
roughness, AO) and triplanar projection samples \emph{each} family on up to
three projections (\texttt{triplanar.glsl:23--29},
\texttt{83--109}, \texttt{119--139}), for up to three materials blended by
barycentric weights (\texttt{solid\_surface.glsl:139--150},
\texttt{191--200}, \texttt{205--215}). Worst case per pixel:

\begin{center}
3 materials $\times$ 4 maps $\times$ 3 projections $=$ \textbf{36 texture fetches},
\end{center}

plus a full tangent-basis construction per projection per material
(\texttt{computeProjectionNormal}, \texttt{triplanar.glsl:71--79}: two
crosses, three normalizes, a normal-map fetch and convention fix-up). The
\texttt{w.x > 0.0} guards skip zero-weight materials, so the common case is
24 fetches (two live materials at triangle boundaries), but single-material
interiors still pay 12 where 4 would do if roughness/AO were not triplanar.
The shipped \texttt{main.frag.spv} contains \textbf{34 static
\texttt{OpImageSample} sites} (measured, \S1) --- the highest fetch-site count
of any non-RT module in the build. Two structural observations:

\begin{itemize}
    \item Roughness and AO are low-frequency signals sampled at full
    triplanar rate. Nothing in the lighting model needs per-projection
    roughness: a single-projection (dominant-axis) sample, or a UV-space
    sample, is visually indistinguishable.
    \item The normal path blends three world-space normals that were each
    built from their own TBN basis, then renormalizes
    (\texttt{triplanar.glsl:107--108}). Blending the \emph{normal-map texels}
    first and building one basis would cut the basis constructions by $3\times$.
\end{itemize}

\textit{Technique: fetch-rate tiering + texel-domain blending.} Sample
roughness/AO once per material on the dominant projection (weight-renormalize
the remaining two for albedo/normal only); blend normal-map texels before the
single TBN transform; consider packing roughness$+$AO$+$height into one
texture so one fetch serves three consumers.

\begin{tcolorbox}[promptbox={Prompt --- C3: collapse the triplanar fetch fan-out}]
Goal: cut worst-case solid fragment fetches from 36 toward 12 with no visible
change. Steps: (1) add \texttt{computeTriplanarRoughnessDominant} /
\texttt{AODominant} helpers sampling only the max-weight projection;
(2) rework \texttt{computeTriplanarNormalUVs} to blend the three normal-map
texels first, then run one \texttt{computeProjectionNormal}-style basis;
(3) keep full-rate paths behind a quality toggle for Maximum. Files:
\texttt{shaders/includes/triplanar.glsl},
\texttt{shaders/includes/solid\_surface.glsl}. Validation: A/B screenshots at
grazing angles (where triplanar matters), static \texttt{OpImageSample} count
in \texttt{main.frag.spv} dropping by at least 40\%, validation-clean.
\\ \textbf{Correction after implementation:} shipped as (1) branchless
dominant-projection roughness/AO (one fetch per material instead of three),
(2) texel-blend-then-transform normals (same texels/weights, one T/B basis
instead of three), (3) zero-weight guards on the non-triplanar albedo/normal/
roughness/AO blocks (weight 0 contributes 0, so skipped fetches are exact).
Measured: static sites 34 $\rightarrow$ 30 ($-12\%$); dynamic worst case
36 $\rightarrow$ 24 ($-33\%$), typical two-material 24 $\rightarrow$ 16,
single-material non-triplanar 12 $\rightarrow$ 4 ($-67\%$). The static-40\%
criterion is replaced: static counting sees each helper once despite
per-material calls, so dynamic multiplicity is the honest metric. Albedo stays
three-projection (look-critical), and no quality toggle was added --- the
deviations are visually null by construction (identical weights and texels;
roughness/AO are low-frequency), verifiable by grazing-angle A/B instead of a
second code path. The superseded full-rate helpers are retained for API
stability; the primary path no longer calls them.
\end{tcolorbox}

\subsection{H4: The Shadow Pass Runs Full Tessellation $+$ Displacement Into Three Cascades, Then Six Fullscreen EVSM Blurs}
\label{sec:h4}

\textbf{Severity: High.}

With shadows enabled (Maximum default), solid geometry is drawn a third and
fourth time: once per cascade cull stream
(\texttt{ShadowRenderer.cpp:801}, \texttt{drawCascadeOnly}) with the full
\texttt{main.vert} $+$ \texttt{main.tesc} $+$ \texttt{main.tese} $+$
\texttt{shadow\_evsm.frag} chain (\texttt{ShadowRenderer.cpp:138--147}), i.e.\
\emph{with} tessellation and \emph{with} displacement height sampling --- the
very work C2 removes from the depth prepass is still paid three times here.
Cascade sizes are 2048/1024/512
(\texttt{ShadowRenderer.cpp:25--26}), and each cascade's moments texture then
pays a horizontal and a vertical 9-tap separable blur
(\texttt{ShadowRenderer.cpp:464--466}, \texttt{evsm\_blur.frag:18,36--38}):
six fullscreen blur passes per frame. Two mitigations already exist and bound
the damage: water is deliberately \emph{not} a shadow caster
(\texttt{ShadowRenderer.cpp} after :801), and Minimal disables shadows
entirely. What remains is the Maximum-default cost: tessellated, displaced
geometry $\times 3$ plus $2048^2 + 1024^2 + 512^2$ blur fragments $\times 2$.

\textit{Technique: cheap shadow-depth geometry + blur gating.} Bind the C2
conservative-depth TES variant (or a no-displacement variant) for the shadow
pipelines --- shadow depth is tolerant to centimetre error, and the existing
depth bias dominates it. Gate the EVSM blur per cascade on screen coverage
(skip the 512-cascade blur when its texel density is below the visible
threshold) and consider a single 5-tap kernel for the two outer cascades.

\begin{tcolorbox}[promptbox={Prompt --- H4: cheap shadow geometry and gated EVSM blur}]
Goal: cut the shadow path to a fraction of its current vertex and blur cost.
Steps: (1) reuse the C2 conservative-depth TES for the shadow pipelines
(\texttt{ShadowRenderer.cpp:138--147}); (2) add per-cascade blur toggles driven
by cascade screen coverage, defaulting the outer cascade to a 5-tap kernel;
(3) keep full quality behind the Maximum preset. Files:
\texttt{vulkan/renderer/ShadowRenderer.cpp},
\texttt{shaders/evsm\_blur.frag}, \texttt{shaders/includes/solid\_tese.glsl}.
Validation: shadow-acne and light-bleed A/B screenshots, timestamp queries on
the cascade and blur passes showing the shadow path shrinking;
validation-clean.
\\ \textbf{Implemented as:} (1) \texttt{SHADOW\_PASS} TES variant
(\texttt{main\_shadow.tese.spv}, 3\,441 vs 7\,881 words): displacement
preserved byte-identically, dead varyings and their interpolation pruned
(interface now 5 inputs + \texttt{POSWORLD}, verified in disassembly); the
C2-style displacement skip was declined for the same hole argument as C2 ---
shadow depths must match main-pass depths within bias. (2) 5-tap EVSM blur
twin for cascades 1--2 (cascade 0 keeps 9 taps), same layout/descriptor
sets/push constants. No quality toggle: outer-cascade filtering differences
are sub-visible by texel density; confirm by shadow A/B (acne/bleed) before
touching the kernel again --- a narrower kernel was already reverted once
for blockiness (\texttt{evsm\_blur.frag} header).
\end{tcolorbox}

\subsection{H5: Every Solid Pixel Pays an RGB-to-HSV-to-RGB Round-Trip With Divergent Branches}
\label{sec:h5}

\textbf{Severity: High.}

The last thing \texttt{shadeSolidSurface} does before writing
\texttt{outColor} is a full color-space round-trip on \emph{every} pixel
(\texttt{solid\_surface.glsl:770--774}):

\begin{lstlisting}[basicstyle=\footnotesize\ttfamily]
vec3 texHSV = rgbToHsv(finalColor);        // 3-way if-chain, mod, divides
texHSV.x = mod(texHSV.x + hsvColor.x, 360.0);
texHSV.y = clamp(texHSV.y * (hsvColor.y * 2.0), 0.0, 1.0);
texHSV.z *= hsvColor.z * 2.0;
finalColor = hsvToRgb(texHSV);             // 6-way if-chain
\end{lstlisting}

\texttt{rgbToHsv} (\texttt{hsv.glsl:24--42}) contains a nested 3-way branch on
the dominant channel plus \texttt{mod}; \texttt{hsvToRgb}
(\texttt{hsv.glsl:4--21}) a 6-way branch. Adjacent pixels routinely take
different paths (channel dominance flips at texture detail), so warps diverge
on top of the ALU. Yet the default vertex tint is identity
(\texttt{math/Vertex.hpp:41}: \texttt{hsv(0.0f, 0.5f, 0.5f)}), for which the
whole round-trip is a no-op (hue $+0$, saturation $\times 1$, value $\times
1$). Paint mode overrides the tint per pixel, but that is already a
uniform-gated branch (\texttt{solid\_surface.glsl:75--91}).

\textit{Technique: hoist the identity test.} Add a single UBO bit
(\texttt{globalHsvIdentity}) set when no painted/brush tint is active and all
resident materials use the default tint; when set, skip both conversions.
Finer: quantize the tint to a $16^3$ LUT indexed in the fragment shader, which
also removes the divergence (one fetch, no branches).

\begin{tcolorbox}[promptbox={Prompt --- H5: gate the per-pixel HSV round-trip}]
Goal: pixels with the identity tint skip both conversions. Steps: (1) add a
UBO flag set by the paint/material system when every active tint is
$(0, 0.5, 0.5)$; (2) wrap \texttt{solid\_surface.glsl:770--774} in it;
(3) optionally replace the conversion pair with a small 3D LUT for
non-identity tints. Files: \texttt{shaders/includes/solid\_surface.glsl},
\texttt{shaders/includes/ubo.glsl}, the UBO upload path. Validation:
pixel-identical screenshots in both states, disassembly showing the
conversions compiled out on the fast path, validation-clean.
\end{tcolorbox}

\subsection{H6: CSM Cascade Selection Burns up to Three Matrix Projections per Pixel --- and RT Hit Shading Repeats It}
\label{sec:h6}

\textbf{Severity: High.}

\texttt{ShadowCalculation} (\texttt{shadows.glsl:32--87}) evaluates
\texttt{fragPosLightSpace.xyz / w} for cascade~0, then up to two full
\texttt{projectToShadowMap} \texttt{mat4} $\times$ \texttt{vec4} multiplies
(:59, :75) for cascades~1--2, plus one or two EVSM fetches with \texttt{exp},
division-heavy Chebyshev math and a light-bleed remap per fetch
(\texttt{evsm.glsl:17--34}). This runs per primary solid pixel \emph{and} per
secondary hit: solid reflection shading calls \texttt{ShadowCalculation} again
for every ray hit (\texttt{solid\_surface.glsl:528--529}), and water RT hits
sample the CSM the same way (report~18 chain). A reflective pixel therefore
pays the cascade-selection prologue two to four times for nearly the same
world position.

\textit{Technique: cascade-index reuse + hit-space amortization.} Pass the
primary pixel's already-computed cascade index (or its projected coordinates)
into the reflection-hit shading path instead of recomputing the selection from
\texttt{worldPos}; hits are centimetres-to-metres from the primary surface and
almost always land in the same cascade. Longer term, replace the three-matrix
ladder with a single atlas transform plus index bits.

\begin{tcolorbox}[promptbox={Prompt --- H6: reuse the primary cascade decision in secondary hits}]
Goal: one cascade selection per pixel, shared by primary and secondary
shading. Steps: (1) refactor \texttt{ShadowCalculation} into
\texttt{selectCascade} (returns index $+$ projected coords) and
\texttt{sampleCascade}; (2) thread the primary result through
\texttt{rtSceneSampleReflectionAlbedo} / \texttt{rtResolveWaterHit} call
sites, reusing it when the hit is within the same cascade bounds and
re-selecting otherwise; (3) verify no double-darkening (combine rule at
\texttt{solid\_surface.glsl:265} unchanged). Files:
\texttt{shaders/includes/shadows.glsl},
\texttt{shaders/includes/solid\_surface.glsl},
\texttt{shaders/includes/rt\_scene\_sample.glsl}. Validation: reflection
screenshots identical, static ALU reduction in the RT fragment variants,
validation-clean.
\end{tcolorbox}

\subsection{H7: A Dead 12-Byte Color Attribute Rides Four Shader Stages in a 64-Byte Vertex}
\label{sec:h7}

\textbf{Severity: High (bandwidth) / Low (effort).}

\texttt{sizeof(Vertex)} is 64~bytes: position 12, \textbf{color 12}, uv 8,
normal 12, brush index 4, pad 4, hsv 12 (\texttt{math/Vertex.hpp:29--39},
\texttt{vulkan/includes/vertex\_layouts.hpp}). The color attribute is fetched
from the merged vertex pool, interpolated by the tessellator, and passed
VS $\rightarrow$ TCS $\rightarrow$ TES $\rightarrow$ FS --- but nothing reads
it: \texttt{main.frag:28} declares \texttt{in vec3 fragColor} and
\texttt{solid\_surface.glsl} never references it (the only \texttt{fragColor}
hit in the solid path is the declaration itself). Dropping the attribute from
the solid pipelines eliminates 12~bytes per vertex of attribute-fetch traffic
and one interpolant across the TCS/TES/FS --- the TES interpolates a dead
varying per generated vertex, which at high tess levels is the dominant
multiplier. Note the subtlety: with \texttt{alignas(16)}, deleting the color
alone leaves \texttt{sizeof(Vertex)} at 64 (52 rounds up), so the stride
itself shrinks only with a repack (e.g.\ fold \texttt{brushIndex} into the
remaining pad, compress \texttt{hsv}) down to 48.

\textit{Technique: delete the attribute, then repack.} Remove
\texttt{ATTR\_COLOR} from the solid pipelines' \texttt{defaultAttributes} (the
\texttt{defaultAttributesFiltered} helper already exists for exactly this),
drop the VS/TCS/TES passthrough, and repack the struct to a 48-byte stride.
If any debug view wants a flat color, use a push constant.

\begin{tcolorbox}[promptbox={Prompt --- H7: remove the dead color attribute}]
Goal: 12~bytes/vertex less attribute traffic, one fewer interpolant, and a
48-byte stride after repacking. Steps: (1) switch the solid pipelines to
\texttt{defaultAttributesFiltered} without \texttt{ATTR\_COLOR}; (2) delete the \texttt{fragColor} passthrough in
\texttt{main.vert}, \texttt{solid\_tesc.glsl}, \texttt{solid\_tese.glsl} and
the \texttt{main.frag:28} declaration; (3) confirm no widget/debug consumer
reads it. Files: \texttt{vulkan/includes/vertex\_layouts.hpp},
\texttt{vulkan/renderer/SolidRenderer.cpp}, \texttt{shaders/main.vert},
\texttt{shaders/main.frag}, \texttt{shaders/includes/solid\_tesc.glsl},
\texttt{shaders/includes/solid\_tese.glsl}. Validation: stride assert
($\mathtt{sizeof(Vertex)}$ or a slimmer solid vertex), pixel-identical
screenshots, vertex-fetch-bytes counters down, validation-clean.
\end{tcolorbox}

\subsection{H8: The Solid TCS Computes Every Edge Level Three Times --- the Water TCS Does It Once}
\label{sec:h8}

\textbf{Severity: High (waste) / trivial (fix).}

\texttt{solid\_tesc.glsl:115--118} computes \texttt{outer0/outer1/outer2} and
writes \texttt{gl\_TessLevelOuter/Inner} unconditionally --- in \emph{all three}
TCS invocations per patch. Each computation is two \texttt{length()} calls,
three material-struct fetches and a smoothstep mix
(\texttt{solid\_tesc.glsl:34--63}), so the patch pays it $3\times$. The water
TCS guards the whole block with \texttt{if (gl\_InvocationID == 0)}
(\texttt{water\_tesc.glsl}), leaving the other two invocations to do pure
passthrough. This is a copy-paste asymmetry, not a design constraint.

\textit{Technique: invocation-0 guard.} Wrap the edge block exactly like the
water TCS does. Alternatively hoist edge-level computation to the CPU per
chunk (levels change slowly), but the guard alone removes two thirds of the
TCS ALU.

\begin{tcolorbox}[promptbox={Prompt --- H8: guard the solid TCS edge block}]
Goal: patch-level work runs once per patch, not three times. Steps:
(1) wrap \texttt{solid\_tesc.glsl:112--129} in
\texttt{if (gl\_InvocationID == 0)} mirroring \texttt{water\_tesc.glsl};
(2) confirm crack-free LOD transitions are unaffected (edge math unchanged,
only placement). Files: \texttt{shaders/includes/solid\_tesc.glsl}.
Validation: tessellation heatmap view identical, TCS instruction count down in
disassembly, validation-clean.
\end{tcolorbox}

% =====================================================================
\section{Medium-Severity Issues}
% =====================================================================

\subsection{M9: The Composite Uses Implicit-LOD Fetches Inside Divergent Flow}
\label{sec:m9}

\textbf{Severity: Medium.}

The \texttt{postprocess.frag} water composite takes the closest of four
\texttt{waterGeomDepthTex} taps with implicit-LOD \texttt{texture()} inside
\texttt{if (waterAlpha > 0.0 \&\& ...)} --- divergent control flow, since
\texttt{waterAlpha} varies per pixel. Implicit derivatives there are
per-quad estimates that go wrong exactly at water silhouettes, the place the
closest-of-four logic exists to protect. The rest of the water path already
learned this lesson: the triplanar helpers carry explicit
\texttt{...Lod} variants precisely because implicit LOD is illegal inside
non-uniform flow (\texttt{triplanar.glsl:31--37}).

\textit{Technique: \texttt{textureLod(..., 0.0)}.} Four-character-class fix;
also hoists cleanly into the H4-style gating (skip the taps where the blur
block already determined there is no water).

\begin{tcolorbox}[promptbox={Prompt --- M9: explicit LOD in the composite depth taps}]
Goal: LOD-correct water-silhouette compositing. Steps: replace the four
\texttt{texture(waterGeomDepthTex, ...)} calls in \texttt{postprocess.frag}
with \texttt{textureLod(..., 0.0)}. Files:
\texttt{shaders/postprocess.frag}. Validation: shoreline A/B screenshots at
pixel zoom, validation-clean (derivative-uniformity warnings gone on strict
drivers).
\end{tcolorbox}

\subsection{M10: \texttt{waterRenderScale} Defaults to 1.0 and Neither Preset Touches It}
\label{sec:m10}

\textbf{Severity: Medium.}

The water offscreen targets honor \texttt{Settings::waterRenderScale}
(\texttt{SceneRenderer.cpp:171--176}, slider $0.25$--$1.0$ in
\texttt{widgets/SettingsWidget.cpp:149}) --- but the default is
\texttt{1.0f} (\texttt{utils/Settings.hpp:143}) and neither
\texttt{applyMaximum} nor \texttt{applyMinimal} sets it
(\texttt{utils/GraphicsSettingsCommand.hpp}). So Minimal pays full-resolution
water geometry, body, column and composite for a surface whose refraction body
is then 9-tap blurred anyway (\S1 composite listing): blur-tolerant work at
full rate by default. Halving the water resolution quarters water fragment,
bandwidth and target-memory cost at a quality loss the blur largely hides.

\textit{Technique: preset-wired render scale.} Set \texttt{waterRenderScale =
0.5} in \texttt{applyMinimal} (keep 1.0 in Maximum), and document the slider
as the override. The composite's closest-of-four depth logic already handles
mismatched target/screen sizes.

\begin{tcolorbox}[promptbox={Prompt --- M10: wire waterRenderScale into the presets}]
Goal: Minimal renders water at half resolution by default. Steps: set
\texttt{settings.waterRenderScale} in \texttt{applyMinimal} ($0.5$) and
\texttt{applyMaximum} ($1.0$); confirm \texttt{SceneRenderer::setWaterRenderScale}
propagation (\texttt{MyApp.cpp:806,1182--1184}) recreates targets correctly.
Files: \texttt{utils/GraphicsSettingsCommand.hpp}. Validation: VRAM and water-pass
timestamps down $\sim$4$\times$ on Minimal, shoreline screenshots acceptable,
validation-clean.
\end{tcolorbox}

\subsection{M11: The Back-Face Pass Re-Dispatches a Full Water Cull Whenever Shadows Are On}
\label{sec:m11}

\textbf{Severity: Medium.}

The water back-face pass (volume thickness) needs culled water commands. With
shadows off it reuses the main water cull output
(\texttt{MyApp.cpp:2469--2473}); with shadows on, the shadow cascade cull
contends for the same buffers, so a second full cull dispatch over the water
slot capacity is recorded every frame (\texttt{MyApp.cpp:2332--2345},
\texttt{taskLocalBackFaceCull}). Cull cost scales with slot \emph{capacity},
not the visible count, so this doubles one of the frame's largest dispatches
in exactly the configuration (Maximum) that is already the most expensive.

\textit{Technique: cull-output sharing across passes.} The cascade cull and
the back-face cull test the same bounds against different frustums; at minimum
give the back-face dispatch its own persistent scratch and compact buffers
sized to the water entry count (not the shared pool), and investigate keying
the back-face draw off the main-pass visible set (back faces are needed only
where front faces survived).

\begin{tcolorbox}[promptbox={Prompt --- M11: stop paying a second full water cull}]
Goal: one water cull dispatch per frame regardless of shadow state. Steps:
(1) scope the task-local back-face cull buffers to the water entry count;
(2) evaluate deriving the back-face draw list from the main water visible set
instead of re-culling; (3) keep the reuse fast path
(\texttt{MyApp.cpp:2469--2473}) as the shadows-off behavior. Files:
\texttt{MyApp.cpp} ($\sim$2332--2345, $\sim$2462--2474),
\texttt{vulkan/renderer/IndirectRenderer.hpp}. Validation: back-face thickness
screenshots identical, one fewer full-capacity dispatch on Maximum,
validation-clean.
\end{tcolorbox}

\subsection{M12: The No-Tess Water Vertex Shader Does up to Six Depth Fetches and Two Inverse-VP Reconstructs per Vertex}
\label{sec:m12}

\textbf{Severity: Medium.}

\texttt{main.vert:87--139} (the \texttt{WATER\_NO\_TESS} path Minimal binds)
samples the solid and back-face depth textures up to six times per vertex and
runs two full inverse-view-projection reconstructions to measure thickness and
shore direction --- vertex-texture-fetch latency plus heavy ALU on every water
vertex, before the fragment stage re-measures depth per pixel anyway
(\texttt{water\_surface.glsl:678}). The vertex measurement seeds only the wave
amplitude fade and foam proximity, both low-frequency signals.

\textit{Technique: chunk-granularity measurement.} Compute per-chunk depth and
shore direction once on the CPU (or in a tiny compute pre-pass) and upload
them as per-draw constants; the VS then does a table lookup instead of six
fetches and two matrix reconstructs. Alternatively quantize the measurement to
a coarse screen-space grid refreshed at half rate.

\begin{tcolorbox}[promptbox={Prompt --- M12: hoist water-vertex depth measurement out of the VS}]
Goal: near-zero texture/ALU cost in the no-tess water VS. Steps: (1) compute
per-chunk water depth $+$ shore direction on the CPU alongside the existing
bounds publication; (2) deliver via the water render UBO or a small storage
buffer indexed by draw id; (3) simplify \texttt{main.vert:87--139} to the
lookup. Files: \texttt{shaders/main.vert},
\texttt{shaders/includes/water\_tese.glsl}, the chunk publication path.
Validation: shoreline foam/wave-fade screenshots identical, VS disassembly
shrinking toward the 682-word tess-path vertex shader, validation-clean.
\end{tcolorbox}

% =====================================================================
\section{Low-Severity / Hygiene Issues}
% =====================================================================

\subsection{L13: Per-Pixel Triplanar \texttt{pow} and Renormalization}
\label{sec:l13}

\texttt{solid\_surface.glsl:110--117} raises the thresholded weights to
\texttt{ubo.triplanarExponent} with \texttt{pow} per pixel, then divides by
the sum. Restrict the exponent to $\{1, 2, 4, 8\}$ (repeated squaring, no
\texttt{pow}) or fold common exponents into a 1D LUT. Minor ALU per pixel,
trivial fix.

\subsection{L14: NaN Diagnostics Execute in Shipping Shaders}
\label{sec:l14}

\texttt{solid\_surface.glsl:156} and \texttt{:182} run
\texttt{isnan}+\texttt{length} checks on every pixel to paint degenerate
normals red. Gate them on \texttt{ubo.debugMode != 0} (or a dedicated
\texttt{NORMAL\_CHECK} build) so the production path pays nothing.

\subsection{L15: Body/Column Targets Are Always Allocated, Even in No-Body Mode}
\label{sec:l15}

\texttt{WaterRenderer::createRenderTargets} unconditionally creates the
full-resolution RGBA16F body and RG16F column targets
(\texttt{WaterRenderer.cpp:398--419}) although report~19 H4 added a
single-attachment pass that never writes them. VRAM-only (no bandwidth when
unused), but at 4K it is $\sim$130\,MB of dead allocation across three
frames. Allocate lazily on first use of the blur-capable variant.

% =====================================================================
\section{Implementation Roadmap and Expected Deltas}
% =====================================================================

\begin{table}[h]
\centering
\small
\begin{tabular}{@{}llc@{}}
\toprule
\textbf{Item} & \textbf{Change} & \textbf{Est.\ frame delta} \\
\midrule
C1 & No-tess depth $+$ shadow (color stays tess.) & $-2$ TCS/TES traversals (depth, 3 shadow) \\
C2 & Prepass gate: opt-in single pass (default on) & $-1$ full traversal when flipped (tess-on regime) \\
C3 & Triplanar fetch tiering (shipped) & $-33$\% worst-case, $-67$\% single-mat \\
H4 & Shadow TES prune $+$ 5-tap outer blur (shipped) & $-40$ ALU/vtx, $-44$\% outer-blur taps \\
H5 & HSV gate & $-$1 round-trip + divergence per solid px \\
H6 & Cascade reuse in secondary hits & $-$1--3 projections per reflective px \\
H7 & Drop dead color attribute & $-12$\,B/vtx traffic, stride 64$\rightarrow$48 \\
H8 & TCS invocation guard & $-67$\% solid TCS ALU \\
M9--M12 & Composite LOD, scale, cull, VS VTF & additive, each 2--10\% \\
L13--L15 & Hygiene & noise-floor / VRAM \\
\bottomrule
\end{tabular}
\end{table}

C1 and C2 compose on the geometry front-end: together they
remove the TCS/TES stages and displacement sampling from the depth prepass
and the three shadow passes, while the color pass keeps the tessellated
family to preserve material blending. C3 is the largest
single fragment-side win. None of the items changes the lighting model, the
CSM authority, or the RT fallback structure, so each can land independently
behind its existing preset toggle.

% =====================================================================
\section{Overall Score}
\label{sec:score}
% =====================================================================

\subsection{Rubric}

Each category is scored 0--10 against a modern explicit-API engine baseline
(dynamic rendering, GPU-driven submission, LOD-graded shading). Weights
reflect frame-time impact for this application's typical scene (dense
tessellated terrain, water, 3-cascade shadows).

\begin{table}[h]
\centering
\small
\begin{tabular}{@{}lcp{7.2cm}@{}}
\toprule
\textbf{Category} & \textbf{Wt.} & \textbf{Score and justification} \\
\midrule
Submission \& culling & 15\% & \textbf{9.0} --- merged GPU cull with stable
slots, indirect-count draws, vegetation/SDF/bbox streams folded into one
dispatch, batched barriers (reports~19--20 work preserved). \\
CPU overhead \& sync & 10\% & \textbf{8.0} --- cached descriptors, single UBO
writes, per-frame mutex discipline; residual per-frame snapshots noted in
report~18 remain. \\
Vertex / geometry & 20\% & \textbf{5.0} --- C1 (always-tessellated solid),
C2 (double geometry), H4 (triple shadow geometry), H7/H8. The front-end is
the weakest subsystem. \\
Fragment fetch efficiency & 20\% & \textbf{5.5} --- C3 (36-fetch worst case),
H5/H6 secondaries; ALU itself is well gated (report~20). \\
Bandwidth \& targets & 15\% & \textbf{6.5} --- RGBA32F water target (report~20
H8) still default-full-res (M10), dead no-body allocations (L15); solid
targets otherwise sane, no MSAA tax. \\
Scalability controls & 10\% & \textbf{7.5} --- presets genuinely disable
subsystems (Minimal), distance LOD, contribution/checkerboard RT gates;
\texttt{waterRenderScale} not preset-wired (M10). \\
Correctness posture & 10\% & \textbf{8.5} --- validation-clean discipline,
feature detection, conservative fallbacks documented in-shader throughout. \\
\bottomrule
\end{tabular}
\end{table}

\subsection{Overall: 6.9/10 --- Good, Held Back by Geometry Multiplicity}

\begin{equation}
0.15(9.0) + 0.10(8.0) + 0.20(5.0) + 0.20(5.5) + 0.15(6.5) + 0.10(7.5) + 0.10(8.5) = \mathbf{6.9}
\end{equation}

The engine's submission architecture (9.0) and correctness posture (8.5) are
its strengths; the score is dragged down by two concentrated weaknesses, both
fixable without architectural changes: geometry is processed two to five
times per frame depending on the pass (5.0), and the solid fragment stage
re-fetches the same texel families instead of once (5.5). Landing C1--C3
alone would move those categories to roughly 7.0/7.5 and the overall score to
\textbf{$\sim$7.8/10}. Nothing in this report suggests a redesign: the fixes
are pipeline variants, shader gates, and attribute deletions behind toggles
that already exist.

\end{document}